\section{Dataset và phân tích nghiệp vụ}

\subsection{Giới thiệu nguồn dữ liệu và cấu trúc dữ liệu}

Nguồn dữ liệu chính được sử dụng trong đề tài là bộ dữ liệu bán lẻ toàn cầu Global Superstore (\texttt{dataset/superstore\_dataset2011-2015.csv}), phản ánh toàn bộ hoạt động đặt hàng và giao nhận của một tập đoàn thương mại điện tử đa quốc gia trong suốt bốn năm tài chính từ ngày 01/01/2011 đến ngày 31/12/2014. Tập dữ liệu gốc bao gồm chính xác 51.290 dòng bản ghi (rows) và 24 cột thuộc tính (columns), với dung lượng lưu trữ vật lý khoảng 12,03 MB và được mã hóa theo bảng mã ký tự Latin-1 (ISO-8859-1 / Windows-1252) do chứa nhiều tên riêng khách hàng và địa danh bằng tiếng Pháp, tiếng Tây Ban Nha và tiếng Đức. Bên cạnh tệp dữ liệu lịch sử chính, hệ thống còn sử dụng tệp dữ liệu gia tăng \texttt{dataset/incremental\_orders\_2015\_q1.csv} gồm 500 bản ghi giao dịch phát sinh trong Quý 1/2015 (có mã định danh \texttt{Row ID} từ 60001 đến 60500) để phục vụ kiểm thử quy trình nạp bổ sung.

Về mặt ngữ nghĩa nghiệp vụ, mỗi dòng bản ghi trong tập dữ liệu đại diện cho một mục hàng chi tiết thuộc một đơn đặt hàng cụ thể (Order Line Item). Một đơn đặt hàng (\texttt{Order ID}) có thể bao gồm một hoặc nhiều dòng bản ghi tương ứng với các mặt hàng khác nhau được khách hàng thanh toán trong cùng một giỏ hàng. Trong quy trình kiểm tra sơ bộ tại tầng tiếp nhận (\texttt{src/etl/bronze/extract.py}), 24 thuộc tính của tập dữ liệu được phân chia thành nhóm 20 thuộc tính bắt buộc (\texttt{REQUIRED\_COLUMNS}) và 4 thuộc tính tùy chọn hoặc bổ trợ (\texttt{OPTIONAL\_COLUMNS}) gồm \texttt{Postal Code}, \texttt{Market}, \texttt{Shipping Cost} và \texttt{Order Priority}. Cấu trúc chi tiết của 24 trường dữ liệu trong nguồn gốc được tổng hợp trong Bảng \ref{tab:raw_schema}.

\begin{table}[H]
\centering
\caption{Cấu trúc các trường thuộc tính trong bộ dữ liệu nguồn Global Superstore}
\label{tab:raw_schema}
\footnotesize
\renewcommand{\arraystretch}{1.22}
\begin{tabularx}{\textwidth}{|l|l|l|X|}
\hline
\textbf{Nhóm nghiệp vụ} & \textbf{Tên thuộc tính gốc} & \textbf{Kiểu dữ liệu đích} & \textbf{Ý nghĩa và đặc điểm dữ liệu} \\
\hline
\multirow{6}{*}{Đơn hàng \& Vận chuyển} 
 & \texttt{Row ID} & \texttt{INTEGER} & Mã định danh duy nhất cho từng dòng chi tiết đơn hàng (1 đến 51290). \\ \cline{2-4}
 & \texttt{Order ID} & \texttt{VARCHAR(50)} & Mã đơn đặt hàng giao dịch (có 25.035 mã đơn hàng phân biệt). \\ \cline{2-4}
 & \texttt{Order Date} & \texttt{DATE} & Ngày khách hàng xác nhận đặt hàng (từ 01/01/2011 đến 31/12/2014). \\ \cline{2-4}
 & \texttt{Ship Date} & \texttt{DATE} & Ngày xuất kho giao hàng thực tế (luôn lớn hơn hoặc bằng \texttt{Order Date}). \\ \cline{2-4}
 & \texttt{Ship Mode} & \texttt{VARCHAR(50)} & Phương thức vận chuyển: Standard Class, Second Class, First Class, Same Day. \\ \cline{2-4}
 & \texttt{Order Priority} & \texttt{VARCHAR(20)} & Mức độ ưu tiên xử lý đơn hàng: Low, Medium, High, Critical. \\
\hline
\multirow{3}{*}{Khách hàng} 
 & \texttt{Customer ID} & \texttt{VARCHAR(50)} & Mã định danh tự nhiên của khách hàng (1.590 mã khách hàng duy nhất). \\ \cline{2-4}
 & \texttt{Customer Name} & \texttt{VARCHAR(100)} & Họ và tên đầy đủ của khách hàng đặt mua. \\ \cline{2-4}
 & \texttt{Segment} & \texttt{VARCHAR(50)} & Phân khúc khách hàng: Consumer (Cá nhân), Corporate (Doanh nghiệp), Home Office. \\
\hline
\multirow{6}{*}{Địa lý \& Thị trường} 
 & \texttt{City} & \texttt{VARCHAR(100)} & Tên thành phố nhận hàng của đơn giao dịch. \\ \cline{2-4}
 & \texttt{State} & \texttt{VARCHAR(100)} & Tên bang, tỉnh hoặc vùng hành chính trực thuộc quốc gia. \\ \cline{2-4}
 & \texttt{Country} & \texttt{VARCHAR(100)} & Tên quốc gia nhận hàng (bao phủ 147 quốc gia trên thế giới). \\ \cline{2-4}
 & \texttt{Postal Code} & \texttt{VARCHAR(20)} & Mã bưu chính khu vực (chỉ có giá trị tại Hoa Kỳ, khuyết thiếu ở các nước khác). \\ \cline{2-4}
 & \texttt{Market} & \texttt{VARCHAR(50)} & Khu vực thị trường vĩ mô: APAC, EU, US, LATAM, EMEA, Africa, Canada. \\ \cline{2-4}
 & \texttt{Region} & \texttt{VARCHAR(50)} & Vùng địa lý chi tiết trực thuộc thị trường (13 vùng kinh doanh). \\
\hline
\multirow{4}{*}{Sản phẩm} 
 & \texttt{Product ID} & \texttt{VARCHAR(50)} & Mã định danh sản phẩm (10.292 mã SKU phân biệt). \\ \cline{2-4}
 & \texttt{Category} & \texttt{VARCHAR(50)} & Ngành hàng chính: Technology, Furniture, Office Supplies. \\ \cline{2-4}
 & \texttt{Sub-Category} & \texttt{VARCHAR(50)} & Tiểu ngành hàng chi tiết (17 phân nhóm sản phẩm con). \\ \cline{2-4}
 & \texttt{Product Name} & \texttt{VARCHAR(255)} & Tên thương mại đầy đủ và mô tả đặc tính của sản phẩm. \\
\hline
\multirow{5}{*}{Độ đo tài chính} 
 & \texttt{Sales} & \texttt{NUMERIC(12,4)} & Doanh thu thuần của dòng sản phẩm (đơn vị: USD, giá trị dương). \\ \cline{2-4}
 & \texttt{Quantity} & \texttt{INTEGER} & Số lượng đơn vị sản phẩm được đặt mua trong dòng giao dịch. \\ \cline{2-4}
 & \texttt{Discount} & \texttt{NUMERIC(5,4)} & Tỷ lệ chiết khấu áp dụng trên dòng hàng (nhận giá trị từ 0,00 đến 0,85). \\ \cline{2-4}
 & \texttt{Profit} & \texttt{NUMERIC(12,4)} & Lợi nhuận ròng thu được (USD, có thể mang giá trị dương hoặc âm). \\ \cline{2-4}
 & \texttt{Shipping Cost} & \texttt{NUMERIC(12,4)} & Chi phí vận chuyển phát sinh cho dòng sản phẩm tương ứng (USD). \\
\hline
\end{tabularx}
\end{table}

\subsection{Exploratory Data Analysis (EDA)}

Quá trình phân tích khám phá dữ liệu (Exploratory Data Analysis -- EDA) được thực hiện toàn diện trên 51.290 bản ghi nhằm làm rõ quy mô hoạt động, xu hướng biến động theo thời gian, cơ cấu đóng góp theo không gian địa lý và hiệu quả sinh lời của từng danh mục sản phẩm. Tính trên toàn bộ giai đoạn 2011--2014, doanh nghiệp đạt tổng doanh thu lũy kế là 12.642.501,91 USD, tổng lợi nhuận ròng đạt 1.467.457,29 USD (tương ứng với biên lợi nhuận gộp toàn cục là 11,61\%), tổng sản lượng hàng hóa bán ra đạt 178.312 đơn vị sản phẩm trên 25.035 đơn đặt hàng, với tỷ lệ chiết khấu trung bình là 14,29\%.

Xét theo chiều thời gian qua bốn năm tài chính, hoạt động kinh doanh của Global Superstore thể hiện xu hướng tăng trưởng mở rộng quy mô rõ rệt và bền vững. Năm 2011, hệ thống ghi nhận 4.440 đơn hàng (8.998 dòng giao dịch) với doanh thu 2.259.450,90 USD và lợi nhuận 248.940,81 USD (biên lợi nhuận 11,02\%). Sang năm 2012, số lượng đơn hàng tăng lên 5.343 đơn (10.962 dòng), mang về 2.677.438,69 USD doanh thu và 307.415,28 USD lợi nhuận (biên lợi nhuận 11,48\%). Đà tăng trưởng tiếp tục bứt phá trong năm 2013 với 6.721 đơn hàng (13.799 dòng), đạt doanh thu 3.405.746,45 USD và lợi nhuận 406.935,23 USD (biên lợi nhuận đạt đỉnh 11,95\%). Đến năm 2014, quy mô giao dịch đạt 8.531 đơn hàng (17.531 dòng), đưa doanh thu lên mức 4.299.865,87 USD và lợi nhuận đạt 504.165,97 USD (biên lợi nhuận 11,73\%). Như vậy, chỉ sau bốn năm vận hành, doanh số hàng năm của doanh nghiệp đã tăng gần gấp đôi (+90,31\% từ năm 2011 đến năm 2014). Khi phân tích sâu theo chuỗi 48 tháng, dữ liệu bộc lộ tính chu kỳ mùa vụ (Seasonality) rất rõ nét: doanh số luôn ở mức thấp nhất vào tháng 1 và tháng 2 đầu năm (dao động từ 91.152 USD đến 98.898 USD trong năm 2011), sau đó tăng tốc mạnh vào tháng 6 và đạt đỉnh cao nhất vào quý 4 hàng năm (đặc biệt là tháng 11 và tháng 12/2014 với doanh thu lần lượt đạt 555.279,03 USD và 503.143,69 USD).

Xét theo chiều không gian địa lý và thị trường tiêu thụ (Bảng \ref{tab:eda_market}), khu vực Châu Á -- Thái Bình Dương (APAC) là thị trường lớn nhất của tập đoàn với 11.002 dòng giao dịch (5.437 đơn hàng), đóng góp 3.585.744,13 USD doanh thu (chiếm 28,36\% toàn cầu) và 436.000,05 USD lợi nhuận (biên lợi nhuận 12,16\%). Đứng thứ hai là thị trường Châu Âu (EU) với 10.000 dòng giao dịch, mang lại 2.938.089,06 USD doanh thu và 372.829,74 USD lợi nhuận, đạt biên lợi nhuận lành mạnh 12,69\% nhờ duy trì mức chiết khấu trung bình thấp (10,31\%). Thị trường Hoa Kỳ (US) và Mỹ Latinh (LATAM) lần lượt đóng góp 2.297.200,86 USD và 2.164.605,17 USD doanh thu. Điểm bất thường nghiêm trọng nhất nằm ở khu vực Trung Đông và Đông Âu (EMEA): mặc dù tạo ra 806.161,31 USD doanh thu từ 5.029 dòng giao dịch, lợi nhuận thu về chỉ vỏn vẹn 43.897,97 USD (biên lợi nhuận tụt xuống mức báo động 5,45\%) do mức chiết khấu trung bình tại khu vực này bị đẩy lên tới 19,61\%. Ngược lại, thị trường Canada tuy có quy mô nhỏ (384 dòng giao dịch, doanh thu 66.928,17 USD) nhưng không áp dụng chiết khấu (0,00\%) nên đạt biên lợi nhuận cao nhất toàn hệ thống là 26,62\%.

\begin{table}[H]
\centering
\caption{Thống kê kết quả kinh doanh theo khu vực thị trường giai đoạn 2011--2014}
\label{tab:eda_market}
\small
\renewcommand{\arraystretch}{1.25}
\begin{tabularx}{\textwidth}{|l|r|r|r|r|r|r|}
\hline
\textbf{Thị trường (Market)} & \textbf{Số đơn hàng} & \textbf{Số dòng (Rows)} & \textbf{Doanh thu (USD)} & \textbf{Lợi nhuận (USD)} & \textbf{Chiết khấu BQ} & \textbf{Biên LN} \\
\hline
APAC & 5.437 & 11.002 & 3.585.744,13 & 436.000,05 & 14,88\% & 12,16\% \\
\hline
EU & 4.593 & 10.000 & 2.938.089,06 & 372.829,74 & 10,31\% & 12,69\% \\
\hline
US & 5.009 & 9.994 & 2.297.200,86 & 286.397,02 & 15,62\% & 12,47\% \\
\hline
LATAM & 5.138 & 10.294 & 2.164.605,17 & 221.643,49 & 13,55\% & 10,24\% \\
\hline
EMEA & 2.462 & 5.029 & 806.161,31 & 43.897,97 & 19,61\% & 5,45\% \\
\hline
Africa & 2.232 & 4.587 & 783.773,21 & 88.871,63 & 15,67\% & 11,34\% \\
\hline
Canada & 201 & 384 & 66.928,17 & 17.817,39 & 0,00\% & 26,62\% \\
\hline
\textbf{Tổng cộng / Toàn cục} & \textbf{25.035} & \textbf{51.290} & \textbf{12.642.501,91} & \textbf{1.467.457,29} & \textbf{14,29\%} & \textbf{11,61\%} \\
\hline
\end{tabularx}
\end{table}

Xét theo phân khúc khách hàng (\texttt{Segment}), nhóm khách hàng Cá nhân (\texttt{Consumer}) chiếm tỷ trọng áp đảo với 818 khách hàng độc lập, thực hiện 13.104 đơn hàng, mang lại 6.507.949,42 USD doanh thu (51,48\% tổng doanh thu) và 749.239,78 USD lợi nhuận (biên lợi nhuận 11,51\%). Phân khúc Doanh nghiệp (\texttt{Corporate}) gồm 476 khách hàng với 7.673 đơn hàng, đóng góp 3.824.697,52 USD doanh thu (30,25\%) và 441.208,33 USD lợi nhuận (biên lợi nhuận 11,54\%). Phân khúc Văn phòng tại gia (\texttt{Home Office}) gồm 296 khách hàng với 4.687 đơn hàng, tạo ra 2.309.854,97 USD doanh thu (18,27\%) và 277.009,18 USD lợi nhuận (đạt biên lợi nhuận cao nhất trong ba phân khúc là 11,99\%).

Xét theo cơ cấu danh mục sản phẩm, ngành hàng Công nghệ (\texttt{Technology}) dẫn đầu về cả doanh thu lẫn khả năng sinh lời với 4.744.557,50 USD doanh thu và 663.778,73 USD lợi nhuận (biên lợi nhuận 13,99\%), nổi bật nhờ hai tiểu ngành hàng chủ lực là Điện thoại (\texttt{Phones} đạt 1.706.824,14 USD doanh thu, 216.717,01 USD lợi nhuận) và Máy sao chụp (\texttt{Copiers} đạt 1.509.436,27 USD doanh thu, 258.567,55 USD lợi nhuận, biên lợi nhuận 17,13\%). Ngành hàng Đồ dùng văn phòng (\texttt{Office Supplies}) tuy có giá trị đơn vị thấp nhưng tần suất mua lại rất lớn (chiếm 31.273 dòng giao dịch), tạo ra 3.787.070,23 USD doanh thu và 518.473,83 USD lợi nhuận (biên lợi nhuận 13,69\%), trong đó nhóm Giấy in (\texttt{Paper}) và Nhãn dán (\texttt{Labels}) đạt biên lợi nhuận lần lượt tới 24,24\% và 20,54\%. Ngược lại, ngành hàng Nội thất (\texttt{Furniture}) tuy đạt doanh thu 4.110.874,19 USD nhưng lợi nhuận chỉ đạt 285.204,72 USD (biên lợi nhuận 6,94\%). Nguyên nhân trực tiếp kéo tụt hiệu quả của ngành Nội thất là tiểu ngành hàng Bàn làm việc (\texttt{Tables}): với 861 dòng giao dịch, nhóm \texttt{Tables} đạt doanh số 757.041,92 USD nhưng gây lỗ ròng tới $-64.083,39$ USD (biên lợi nhuận âm $-8,46\%$) do mức chiết khấu bình quân lên đến 29,07\% kết hợp cùng chi phí vận chuyển hàng cồng kềnh rất cao.

\subsection{Đánh giá chất lượng và rủi ro dữ liệu}

Thông qua quá trình kiểm tra chất lượng dữ liệu nguồn (Data Profiling), nhóm thực hiện phát hiện năm vấn đề kỹ thuật và rủi ro dữ liệu cần phải được xử lý triệt để trong quy trình ETL trước khi đưa vào mô hình phân tích đa chiều:

Vấn đề thứ nhất là sự bất đồng nhất nghiêm trọng về định dạng ngày tháng tại hai trường \texttt{Order Date} và \texttt{Ship Date}. Trong cùng một tệp CSV, các bản ghi từ ngày 01 đến ngày 12 hàng tháng thường được lưu theo dạng \texttt{D/M/YYYY} sử dụng dấu gạch chéo (ví dụ: \texttt{1/1/2011}), trong khi các bản ghi từ ngày 13 trở đi lại được lưu theo định dạng \texttt{DD-MM-YYYY} sử dụng dấu gạch ngang (ví dụ: \texttt{13-01-2011}). Nếu sử dụng bộ phân tích ngày tháng mặc định của hệ quản trị cơ sở dữ liệu mà không phân nhánh biểu thức chính quy (Regular Expression), quá trình chuyển đổi kiểu dữ liệu sẽ phát sinh lỗi \texttt{DateTimeParseException} hoặc đảo ngược ngày với tháng làm sai lệch hoàn toàn tính mùa vụ của chuỗi thời gian.

Vấn đề thứ hai là tỷ lệ khuyết thiếu dữ liệu (Missing Values) rất lớn tại thuộc tính \texttt{Postal Code}. Trong tổng số 51.290 dòng giao dịch, chỉ có 9.994 dòng thuộc lãnh thổ Hoa Kỳ là có mã bưu chính hợp lệ (tương ứng với 632 tổ hợp mã bưu chính duy nhất được trích xuất tại tầng Bronze), trong khi 41.296 dòng còn lại (chiếm 80,51\% toàn tập dữ liệu) mang giá trị rỗng (\texttt{NULL} hoặc chuỗi trống). Nếu sử dụng trực tiếp cột \texttt{Postal Code} chứa giá trị \texttt{NULL} làm một phần của khóa tự nhiên tổ hợp trong bảng chiều địa lý (\texttt{dim\_geography}), phép toán \texttt{UNIQUE} và phép kết nối \texttt{JOIN} chuẩn SQL sẽ thất bại do trong đại số quan hệ, biểu thức \texttt{NULL = NULL} không trả về giá trị đúng (\texttt{UNKNOWN}). Do đó, việc chuẩn hóa các giá trị khuyết thiếu về một nhãn mặc định tường minh (\texttt{'Unknown'}) là yêu cầu bắt buộc.

Vấn đề thứ ba là hiện tượng đa trị và phụ thuộc hàm không chặt chẽ trong phân cấp địa lý và danh mục sản phẩm. Cụ thể, một số thành phố có cùng tên gọi và cùng quốc gia nhưng lại xuất hiện ở các khu vực quản lý (\texttt{Market} / \texttt{Region}) khác nhau trong dữ liệu lịch sử (ví dụ các đơn hàng tại Vienna, Áo vừa xuất hiện dưới thị trường \texttt{EU} vừa xuất hiện dưới thị trường \texttt{EMEA}), hoặc cùng một mã định danh sản phẩm tự nhiên (\texttt{Product ID}) nhưng có sự khác biệt nhỏ về chuỗi tên sản phẩm (\texttt{Product Name}) giữa các năm. Nếu thiết kế ràng buộc duy nhất không đầy đủ hoặc không sử dụng kỹ thuật chọn bản ghi đại diện mới nhất (\texttt{SELECT DISTINCT ON ... ORDER BY order\_date DESC}), lệnh nạp chiều bằng \texttt{ON CONFLICT DO UPDATE} trong PostgreSQL sẽ lập tức dừng thực thi với lỗi \texttt{CardinalityViolation: ON CONFLICT DO UPDATE command cannot affect row a second time}.

Vấn đề thứ tư là rủi ro trùng lặp lô nạp (Duplicate Batch Ingestion) và ép kiểu số học. Tất cả 24 cột khi đọc từ tệp văn bản CSV vào tầng thô đều ở dạng chuỗi ký tự (\texttt{TEXT}) có thể kèm theo khoảng trắng thừa ở hai đầu. Nếu người vận hành vô tình kích hoạt quy trình nạp nhiều lần trên cùng một tệp nguồn mà không có cơ chế kiểm tra chữ ký số của tệp và cơ chế chống ghi đè tại khóa chính \texttt{row\_id}, các chỉ số cộng được như doanh thu và lợi nhuận sẽ bị nhân bội sai thực tế.

\subsection{Xác định yêu cầu và câu hỏi nghiệp vụ}

Từ kết quả phân tích khám phá và các vấn đề thực tiễn nêu trên, hệ thống Kho dữ liệu và Hỗ trợ quyết định được đặt ra yêu cầu phải chuẩn hóa 13 chỉ số đo lường hiệu quả cốt lõi (KPIs) thuộc ba nhóm tài chính, sản phẩm và khách hàng, được định nghĩa toán học chặt chẽ trong Bảng \ref{tab:kpi_definitions}.

\begin{table}[H]
\centering
\caption{Hệ thống các chỉ số đo lường hiệu quả kinh doanh (KPIs) chuẩn hóa trong hệ thống}
\label{tab:kpi_definitions}
\footnotesize
\renewcommand{\arraystretch}{1.3}
\begin{tabularx}{\textwidth}{|l|l|l|X|}
\hline
\textbf{Nhóm KPI} & \textbf{Tên chỉ số} & \textbf{Công thức xác định} & \textbf{Phân loại \& Ý nghĩa quản trị} \\
\hline
\multirow{6}{*}{Tài chính} 
 & Tổng doanh thu ($S$) & $S = \sum \text{sales}$ & Độ đo cộng được (Additive); phản ánh tổng quy mô doanh số thuần. \\ \cline{2-4}
 & Tổng lợi nhuận ($P$) & $P = \sum \text{profit}$ & Độ đo cộng được; đo lường hiệu quả sinh lời thực tế sau chi phí. \\ \cline{2-4}
 & Biên lợi nhuận ($M$) & $M = \frac{\sum \text{profit}}{\sum \text{sales}} \times 100\%$ & Độ đo không cộng được (Non-additive); đánh giá tỷ suất sinh lời trên doanh thu. \\ \cline{2-4}
 & Chiết khấu bình quân ($\bar{D}$) & $\bar{D} = \frac{1}{N}\sum \text{discount} \times 100\%$ & Độ đo bán cộng (Semi-additive); giám sát mức độ giảm giá trung bình. \\ \cline{2-4}
 & Giá trị đơn trung bình ($AOV$) & $AOV = \frac{\sum \text{sales}}{\text{CountDistinct}(\text{order\_id})}$ & Non-additive; đo lường quy mô chi tiêu bình quân trên mỗi giỏ hàng ($504,99$ USD). \\ \cline{2-4}
 & Tăng trưởng tháng ($g_{\text{MoM}}$) & $g_{\text{MoM}} = \frac{S_t - S_{t-1}}{S_{t-1}} \times 100\%$ & Non-additive; theo dõi động lượng tăng trưởng doanh thu giữa hai tháng liên tiếp. \\
\hline
\multirow{3}{*}{Sản phẩm} 
 & Sản lượng tiêu thụ ($Q$) & $Q = \sum \text{quantity}$ & Additive; tổng số đơn vị hàng hóa xuất kho (178.312 đơn vị). \\ \cline{2-4}
 & Tỷ trọng tích lũy Pareto & $P_{\text{cum}}(k) = \frac{\sum_{i=1}^k S_i}{\sum_{i=1}^M S_i} \times 100\%$ & Non-additive; xác định nhóm 20\% danh mục mang lại 80\% doanh thu/lợi nhuận. \\ \cline{2-4}
 & Độ nâng kết hợp ($\text{Lift}$) & $\text{Lift}(A \rightarrow B) = \frac{P(A \cap B)}{P(A) \cdot P(B)}$ & Đánh giá mức độ phụ thuộc mua kèm giữa hai sản phẩm để đề xuất bán chéo. \\
\hline
\multirow{3}{*}{Khách hàng} 
 & Độ mới giao dịch (\texttt{Recency}) & $R = \text{Date}_{\text{ref}} - \max(\text{order\_date})$ & Số ngày kể từ lần mua hàng gần nhất của khách hàng tính đến mốc tham chiếu. \\ \cline{2-4}
 & Tần suất mua (\texttt{Frequency}) & $F = \text{CountDistinct}(\text{order\_id})$ & Tổng số đơn đặt hàng độc lập mà một khách hàng đã thực hiện. \\ \cline{2-4}
 & Giá trị chi tiêu (\texttt{Monetary}) & $M_{\text{cust}} = \sum \text{sales}$ & Tổng giá trị doanh thu mà khách hàng đóng góp trong toàn bộ vòng đời. \\
\hline
\end{tabularx}
\end{table}

Dựa trên hệ thống KPI này, kho dữ liệu và hệ hỗ trợ quyết định được thiết kế để trả lời trực tiếp bốn nhóm câu hỏi nghiệp vụ trọng tâm. Nhóm câu hỏi thứ nhất về hiệu quả doanh thu và địa lý: Xu hướng doanh thu và lợi nhuận biến thiên như thế nào theo từng tháng, quý, năm tại từng thị trường và quốc gia, đồng thời khu vực nào đang rơi vào trạng thái doanh số cao nhưng biên lợi nhuận âm do lạm dụng chiết khấu? Nhóm câu hỏi thứ hai về quản trị danh mục sản phẩm: Những tiểu ngành hàng và mã sản phẩm SKU cụ thể nào nằm trong nhóm đứng đầu về lợi nhuận (Top Performers) hoặc nhóm gây thất thoát dòng tiền lớn nhất (Loss Makers) cần phải cắt giảm chiết khấu? Nhóm câu hỏi thứ ba về hành vi khách hàng và giỏ hàng: Cơ cấu đóng góp doanh thu theo các phân khúc Consumer, Corporate, Home Office ra sao và những cặp sản phẩm/tiểu ngành hàng nào có xác suất mua kèm cao nhất để thiết kế gói khuyến mãi? Nhóm câu hỏi thứ tư về dự báo và mô phỏng chính sách (DSS): Doanh số dự kiến trong 6 tháng tiếp theo sẽ đạt mức bao nhiêu trong khoảng tin cậy 95\%, và nếu ban giám đốc điều chỉnh tăng hoặc giảm tỷ lệ chiết khấu một mức $\Delta d$ tại một thị trường/ngành hàng cụ thể thì lợi nhuận ròng sẽ thay đổi như thế nào để ra quyết định phê duyệt hay từ chối chính sách?
